\section{Related Work}\label{sec:background}

\paragraph{Auctions.}
There is a vast quantity of theoretical and experimental literature on auctions. While \citet{krishna2009auction}'s textbook and \citet{kagel2020handbook}'s handbook provide invaluable general resources, we focus on citations relevant to the results in this paper.
Starting with the IPV case, the benchmark of revenue equivalence in the risk-neutral case is developed in \citet{myerson1981optimal}. The experimental evidence for departures due to risk-aversion in the FPSB auction is documented in \citet{coppinger1980incentives} and \citet{cox1988theory}'s survey. The experimental evidence for the common error of bidding above one's value in the SPSB auction is documented in \citet{kagel1993independent}; as we show in Section~\ref{sec:benchmark_experimental}, LLM bidders instead tend to bid \emph{below} value in the SPSB auction, an important qualitative divergence from human behavior.
For common value settings, experimental evidence for the winner's curse was first given by \citet{bazerman1983won}; we primarily follow the experimental design and analysis framework in \citet{kagel1986winner} and \citet{kagel1987information}. For eBay-style auctions, the phenomenon of bid sniping has been studied  by \citet{roth2002last} and \citet{ockenfels2006late}, while \citet{einav2011learning} provide an empirical study of hidden reserves on eBay. 

\paragraph{Mechanism descriptions and framing.} There is work that explores how the description or framing of a mechanism affects behavior. \citet{gonczarowski2023strategyproofness} present a method to highlight strategy-proofness with short, one-sentence descriptions in an extensive-form menu, though their menu-description treatment of an SPSB auction did not perform well empirically. \citet{breitmoser2022obviousness} show that re-framing a static sealed-bid auction as an ascending clock auction can improve the rate of truth-telling behavior.

\paragraph{Automated and computational mechanism design.}
Our work also connects to computational and automated mechanism design, which uses algorithms and optimization to search over (or learn) allocation and payment rules in complex environments.
A prominent learning-based line trains neural networks to maximize an objective (e.g., revenue) while controlling incentive-compatibility violations, as in RegretNet and follow-ups \citep{dutting2024optimal,curry102automated}. More recent architectures tighten incentive guarantees by building strategyproofness into the mechanism class itself rather than relying only on soft regret penalties, including GemNet \citep{wang2024gemnet}. These ideas have been extended to richer design problems such as auctions with budgets \citep{feng2018deep}, combinatorial auctions \citep{wang2025bundleflow}, and market-design settings beyond standard auctions such as data markets \citep{ravindranath2023data}.
Our contribution is orthogonal: rather than learning a new mechanism, we provide a benchmarking and reconstruction methodology for comparing \emph{behavioral bid distributions}, LLM-generated versus human, across canonical formats, which can be used as an input to (and diagnostic for) automated design pipelines.

\paragraph{LLMs as simulated agents.} While LLMs  are able to generate human-like text and reasoning \citep{achiam2023gpt, bubeck2023sparks}, results also show that they can  have limited planning abilities and
exhibit various cognitive biases \citep{wan2023kelly}. Contemporaneous with and subsequent to an early version of our paper, a growing literature has explored using these human-like AI models as simulated agents in economics and social science~\citep{aher2023using, park2023generative, horton2023large, manning2024automated}. Our work fits squarely into this agenda: we use LLMs as synthetic experimental subjects to probe whether they can serve as proxies for human behavior in economic mechanisms.

\paragraph{LLMs in auctions.} Since the first version of this paper, there has been a proliferation of work on LLMs in auction settings, which can be broadly divided into three categories.

\emph{LLMs as first-class participants.} The first category treats LLMs as first-class entities---bidders or agents who actively participate in a marketplace---and studies their strategic behavior in its own right. \citet{fish2024algorithmic} study collusive behavior in first-price sealed-bid auctions. \citet{chen2023put} study how to make LLMs better at playing auctions. More recent work includes studies of information disclosure with LLM agents \citep{infobid2025}, comparisons across frontier models in various auction formats \citep{lu2025claude}, and early feasibility checks on LLM bidding behavior \citep{tsuchihashi2023homo}. \citet{auctionnet2024} also study a budget-constrained repeated ad auction (GSP-style), where bidders optimize campaign objectives over time. Much of this work uses auctions primarily as agent benchmarks for planning or collusion, without validating distributional similarity to human lab bidding, or focuses on narrow mechanism families without cross-format replication.

\emph{LLMs as proxies for human behavior in economic design.} The second category---into which our paper falls---asks whether LLM agents can serve as useful proxies for human participants, with the goal of informing mechanism design. \citet{corrigan2025usefulness} consider distributional similarity, showing that a light-touch calibration step using a small sample of human bids can bring LLM bid distributions closer to human benchmarks. Their approach is complementary to ours: while calibration is a promising direction for improving LLM fidelity in specific settings, our focus is on out-of-the-box benchmarking across multiple canonical formats, which is necessary to understand the baseline reliability of LLM proxies before calibration is applied.

\emph{Auction design for LLM-powered platforms.} A third, distinct line of work studies the design of auction mechanisms \emph{within} LLM-powered platforms---for example, how to allocate and price advertising slots or sponsored content when an LLM generates the display surface \citep{duetting2024mechanism, ramage2024adauctions}. Here the research question is about mechanism design for a new marketplace, rather than about using LLMs to simulate bidder behavior, and is thus orthogonal to our agenda.

\paragraph{LLMs and game theory.}
Beyond auctions, there is growing interest in connecting LLMs with game-theoretic modeling of strategic environments. Some work uses LLMs to help represent or extract formal games from unstructured descriptions—for example, deriving extensive-form structure, incentives, or interaction graphs from narratives \citep{daskalakis2024charting}. Other work integrates LLMs into broader mechanism-design pipelines or studies strategic interaction with LLM agents in general-sum settings \citep{duetting2024mechanism,huang2025accelerated,soumalias2025llm, jiang2025incentive, wang2026llm}. 

\bigskip


In this work, we treat LLMs  as synthetic experimental subjects and benchmark their bid distributions against decades of human laboratory evidence across different  auction environments, using a single distributional metric (SMAD) and a moment-matching reconstruction procedure that makes legacy experimental results directly comparable.
